% VI/20 challenge solutions -- full-solutions edition only.

\paragraph{VI.20.C1 solution (Four-chart cocycle audit).}
% VI20_FULL_CHALLENGE_SOLUTION C01
\begin{solution}
A concrete model is obtained by taking four copies $X_i=\Spec k[t]$ and using the common principal open $D(t)$ as every nontrivial overlap.  Choose a spanning tree of the chart graph, for example the transitions
\[
\varphi_{12},\qquad \varphi_{23},\qquad \varphi_{34},
\]
and take each to be an automorphism of $D(t)=\Spec k[t,t^{-1}]$.  For definiteness, let them all be the identity.  The inverse law forces $\varphi_{21},\varphi_{32},\varphi_{43}$.  The cocycle then forces
\[
\varphi_{13}=\varphi_{23}\circ\varphi_{12},\quad
\varphi_{24}=\varphi_{34}\circ\varphi_{23},\quad
\varphi_{14}=\varphi_{34}\circ\varphi_{23}\circ\varphi_{12},
\]
and the reverse transitions are their inverses.

More generally, the three tree transitions may be arbitrary chosen automorphisms of the common overlap; all other transitions are then defined by composing along the unique tree path.  To check coherence, take any three indices $i,j,k$.  By construction, the transition from $i$ to $k$ is the composition along the path from $i$ to $k$.  Concatenating the paths from $i$ to $j$ and $j$ to $k$, and cancelling immediate inverse edges, gives the same path from $i$ to $k$.  Hence
\[
\varphi_{ik}=\varphi_{jk}\circ\varphi_{ij}
\]
on the common overlap.  Thus inverse and cocycle laws hold, so the four charts form coherent gluing data.  The example also shows the efficient principle: on a connected overlap graph, transitions on a spanning tree are free data; the cocycle forces the remaining transitions whenever all required triple overlaps are present.
\end{solution}

\paragraph{VI.20.C2 solution (Doubled-origin morphisms).}
% VI20_FULL_CHALLENGE_SOLUTION C02
\begin{solution}
Let $X$ be the doubled-origin line, obtained from two copies $U_1=U_2=\Spec k[t]$ glued by the identity on $D(t)$.  A morphism $X\to\A^1_k=\Spec k[z]$ is, by the gluing mapping property, a pair of morphisms $U_i\to\A^1_k$ agreeing on the overlap.  By affine anti-equivalence these chart morphisms correspond to polynomials
\[
p_1(t),p_2(t)\in k[t]
\]
via $z\mapsto p_i(t)$.  Agreement on $D(t)$ means
\[
p_1(t)=p_2(t)\quad\text{in }k[t,t^{-1}].
\]
The localization map $k[t]\hookrightarrow k[t,t^{-1}]$ is injective, so $p_1=p_2$ already in $k[t]$.

Thus every morphism is determined by one polynomial $p(t)$, used on both charts.  Therefore
\[
\operatorname{Hom}_k(X,\A^1_k)\cong k[t].
\]
This agrees exactly with the global-section computation
\[
\Gamma(X,\OO_X)\cong k[t],
\]
as predicted by the universal property of the affine line.  The example is instructive: although global functions do not make $X$ affine, they still classify maps from $X$ to an affine line.
\end{solution}

\paragraph{VI.20.C3 solution (Atlas comparison).}
% VI20_FULL_CHALLENGE_SOLUTION C03
\begin{solution}
Write the two affine covers as
\[
X=\bigcup_{i=1}^r U_i=\bigcup_{j=1}^s V_j,\qquad U_i=\Spec A_i,\quad V_j=\Spec B_j.
\]
For each pair $(i,j)$ and each point $x\in U_i\cap V_j$, VI.20.P2 supplies an open neighborhood $W_{ij,x}$ satisfying
\[
W_{ij,x}=D(f_{ij,x})\subseteq U_i,
\qquad
W_{ij,x}=D(g_{ij,x})\subseteq V_j.
\]
The collection of all such $W_{ij,x}$ covers $X$ and therefore is a common affine refinement of both original atlases.

The coordinate ring of each refined chart is simultaneously
\[
\Gamma(W_{ij,x},\OO_X)
\cong(A_i)_{f_{ij,x}}
\cong(B_j)_{g_{ij,x}}.
\]
Hence transition maps between the two original coordinate systems are localization isomorphisms.  If two refined charts overlap, apply P2 once more to their intersection; after this secondary distinguished refinement, every transition between refined charts is again represented by an isomorphism between localizations.  Since all these transitions arise from identity maps on literal open subsets of $X$, their restrictions automatically satisfy the cocycle law.  This yields an explicit common distinguished atlas with entirely algebraic transition data.
\end{solution}

\paragraph{VI.20.C4 solution (Non-affine two-chart experiment).}
% VI20_FULL_CHALLENGE_SOLUTION C04
\begin{solution}
Consider three representative two-chart gluings.

First, glue two copies of $\A^1_k$ along the whole chart by the identity.  The equalizer is $k[t]$, and the gluing is actually $\A^1_k$, because full-overlap gluing removes the duplicate chart.  This example is affine.

Second, glue two copies of $\A^1_k$ along $D(t)$ by the identity.  The equalizer is again
\[
\{(p,q):p=q\text{ in }k[t,t^{-1}]\}\cong k[t].
\]
However, the quotient has two distinct origins.  If it were affine, reconstruction from its global ring would identify it with $\Spec k[t]$, which has only one origin.  Therefore the doubled-origin line is not affine.

Third, glue $\Spec k[t]$ and $\Spec k[u]$ on $D(t),D(u)$ by $u=t^{-1}$.  The equalizer is
\[
k[t]\cap k[t^{-1}]=k
\]
inside $k[t,t^{-1}]$.  If this scheme were affine, it would be $\Spec k$, a one-point space, whereas the glued scheme contains at least the two chart origins.  Hence it is not affine.

These examples show how the equalizer is used.  If a two-chart gluing is assumed affine, then the canonical map to the spectrum of the equalizer must be an isomorphism.  A mismatch of underlying topology (or local rings) proves non-affineness.  Conversely, matching global rings alone does not prove affineness, as the doubled-origin example demonstrates.
\end{solution}

\paragraph{VI.20.C5 solution (Gluing theorem from mapping property).}
% VI20_FULL_CHALLENGE_SOLUTION C05
\begin{solution}
Start with the quotient topological space $X$ and the compatible-family sheaf.  Let $Y$ be a scheme.  If $f:X\to Y$ is a morphism, restriction to the open chart images produces morphisms $f_i:X_i\to Y$ satisfying $f_i|_{U_{ij}}=f_j\circ\varphi_{ij}$.  Conversely, from such a compatible family define $f([x])=f_i(x)$.  Compatibility makes this well defined, the final topology makes it continuous, and the sheaf pullbacks glue because they agree on overlaps.  Locality of stalk maps is checked in a chart.  This proves
\[
\operatorname{Hom}(X,Y)
\cong
\{(f_i): f_i|_{U_{ij}}=f_j\circ\varphi_{ij}\}.
\]

Now suppose $X'$ is another gluing of the same data.  The chart embeddings $X_i\to X'$ form a compatible family, so the mapping property gives a unique $F:X\to X'$.  Reversing the roles gives $G:X'\to X$.  Both composites restrict to the identity on every chart, hence uniqueness forces $GF=\operatorname{id}_X$ and $FG=\operatorname{id}_{X'}$.  This rederives uniqueness of the gluing.

For an existing scheme covered by affine opens $U_i$, use the identity transitions on $U_i\cap U_j$.  The inclusions $U_i\hookrightarrow X$ are compatible, so the mapping property gives a morphism from the reconstructed gluing to $X$; it is an isomorphism on every chart and therefore globally.  This rederives reconstruction from affine covers.

Finally, compare one-step and staged gluing of finitely many charts.  Both constructions contain the same chart morphisms with the same transition compatibilities.  Applying the mapping property in both directions yields inverse chart-compatible morphisms.  Hence staged gluing is associative up to unique canonical isomorphism.
\end{solution}
